\chapter{Conclusion}

Ce travail avait pour objectif la conception intégrale d'un lecteur audio portable, compact et fonctionnel, du schéma électrique au routage du circuit imprimé, puis au firmware embarqué et aux mesures de validation. Cet objectif est atteint. L'appareil est assemblé, fonctionne sur batterie et a été caractérisé sur banc de mesure.

Le bilan face au cahier des charges est clair. Dix des onze exigences minimales sont pleinement satisfaites. La onzième, l'autonomie, l'est dans les conditions établies à la section 6.7. Aucune fonction imposée n'a été abandonnée en cours de projet. La majorité des objectifs supplémentaires d'interface ont également été livrés, et dix fonctions absentes du cahier des charges, orientées vers la robustesse et le diagnostic, ont été ajoutées.

Toutes les étapes fixées en introduction ont été franchies, des spécifications au routage, puis à l'assemblage, au firmware et à la campagne de mesure, la carte quatre couches étant fonctionnelle dès son premier lancement de fabrication. Les mesures confirment la qualité de l'intégration. La sortie filaire présente un THD+N de $-94{,}7$\,dB à 1\,kHz, au niveau de la valeur typique du convertisseur, ce qui montre que l'étage analogique n'ajoute pas de distorsion mesurable malgré la présence d'une alimentation à découpage et d'un émetteur Bluetooth sur la même carte. Le rapport signal sur bruit de l'appareil livré s'établit à 106,2\,dB. La liaison Bluetooth reste stable jusqu'à une vingtaine de mètres, soit le double de l'exigence. Le firmware assure une lecture continue jusqu'à 24 bits et 192\,kHz tout en pilotant l'affichage et la navigation. L'ensemble tient dans un boîtier de 112,5\,$\times$\,69,5\,$\times$\,19,8\,mm pour une masse de [Insérer masse ici]\,g, conforme au format de poche exigé.

Ces résultats s'accompagnent de limites clairement identifiées. L'autonomie mesurée atteint 13\,h\,28 en veille, 9\,h\,14 en lecture filaire et 5\,h\,48 en Bluetooth. L'objectif de dix heures n'est donc pas atteint en lecture. Il est manqué de peu sur la sortie filaire et plus nettement en Bluetooth. La mise au point a par ailleurs révélé un défaut de conception sur le réseau d'écrêtage du bouton d'allumage, dont l'analyse a abouti au montage corrigé présenté au chapitre 7. Enfin, l'absence d'horloge temps réel sur la carte prive l'appareil de trois objectifs supplémentaires liés à l'affichage de l'heure. La solution retenue pour une révision, une horloge sauvegardée par pile bouton, est établie au même chapitre.

Sur le plan personnel, ce projet m'a beaucoup appris sur la conception de bout en bout d'un produit qui vise le grand public. Le choix des composants s'est révélé plus exigeant que prévu, chaque candidat devant être évalué sur ses performances mais aussi sur sa disponibilité réelle en stock. L'intégration du firmware avec des blocs matériels aussi différents a constitué une autre difficulté majeure, chaque bloc imposant ses propres contraintes de communication et de temps. Enfin, mener l'appareil jusqu'à son état final m'a montré que délivrer un produit de qualité demande autant de rigueur dans la vérification que dans la conception elle-même.

Les perspectives futures sont détaillées au chapitre 7 et aucune ne remet en cause l'architecture générale de l'appareil. Parmi elles, l'ajout d'une horloge temps réel et d'une détection de présence du jack apporterait le plus à l'usage quotidien pour un effort mesuré. Le prototype constitue ainsi une base solide pour une seconde révision.



% La signature est optionnelle et pas forcément nécessaire...
\vfil
\hspace{8cm}\makeatletter\@author\makeatother\par
\hspace{8cm}\begin{minipage}{5cm}
    % Place pour signature numérique
    \printsignature
\end{minipage}